\section{问题二：带安全裕度的日前计划模型}

\subsection{因果预测}

题目没有给出问题二可用的日前负载与光伏预测，只给了 2025 年全年的实际值。为保证日前决策的因果性，本文取\textbf{此前 7 天的历史均值}作为当天逐区间的预测：
\begin{equation}
L_t^{\mathrm{fc}}=\frac1{7}\sum_{k=1}^{7}L_t^{(d-k)},\qquad
PV_t^{\mathrm{fc}}=\frac1{7}\sum_{k=1}^{7}PV_t^{(d-k)},
\label{eq:fc}
\end{equation}
其中 $L_t^{(d-k)}$ 表示第 $d-k$ 天第 $t$ 个区间的实际值。第 $d$ 天的决策只用到第 $d-1$ 天及更早的数据，满足因果约束。2025 年 1 月为滚动起步的预热月，其结果不计入 2 月至 12 月的交付范围。

在 2025 年 2 月至 12 月上，该预测的精度为：负载平均绝对误差 $130.2\,\mathrm{kWh}$/区间，光伏平均绝对误差 $25.6\,\mathrm{kWh}$/区间。负载误差明显大于光伏误差，因此净负荷误差 $e_t$ 主要由负载主导。

\subsection{计划模型}

第 $d$ 天 0:00 在预测曲线上求解与问题一结构相同的线性规划，但约束改为日循环而非固定首尾电量：
\begin{equation}
s_{-1}=s_{N-1}=s^{(d)}_0,\qquad
S^{\min}\le s_t\le S^{\max},
\label{eq:p2soc}
\end{equation}
其中 $s^{(d)}_0$ 是当天起始的实际储电量。目标函数仍为 $\min\sum_t p_t x_t$。

\subsection{安全裕度：报童模型}

按预测原样购电并不是最优的。由式~\eqref{eq:identity}，执行缺额等于 $e_t-w_t$，而 5 倍的紧急购电价使得“少买”远比“多买”昂贵。设区间 $t$ 上在计划中额外多买 $y_t\ \mathrm{kWh}$（即在计划内多弃光 $y_t$）：

\begin{itemize}
  \item 若实际并不缺电，多买的 $y_t$ 只能弃掉，边际成本为 $p_t$；
  \item 若实际缺电 $e_t>y_t$，多买可以少付 $(e_t-y_t)$ 部分的 5 倍紧急购电费，边际收益为 $5p_t\cdot P(e_t>y_t)$。
\end{itemize}

令两者相等，得到经典报童模型\cite{porteus2002,zipkin2000}的最优条件：
\begin{equation}
p_t=5p_t\,P(e_t>y_t)
\quad\Longrightarrow\quad
P(e_t>y_t)=\frac{1}{5}
\quad\Longrightarrow\quad
y_t=F_{e_t}^{-1}(0.8),
\label{eq:newsvendor}
\end{equation}
即最优安全裕度就是预测误差分布的 $\mathbf{0.8}$ 分位点。实际计算时用此前 28 天同一区间的误差样本估计该分位点，并截断到非负。由于 $y_t$ 只影响购电量与弃光量、不改变储能充放电决策，可以把裕度直接叠加在规划解之上：
\begin{equation}
x_t \leftarrow x_t^{\text{LP}}+y_t,\qquad w_t\leftarrow w_t^{\text{LP}}+y_t .
\label{eq:margin}
\end{equation}

\subsection{求解结果}

表 4 给出不同分位点下的全年费用。实际最优分位点恰为 $q=0.80$，与式~\eqref{eq:newsvendor} 的理论值一致；相比不加裕度，全年节省 $693$ 万元（$29\%$）。

\threelinetable{表 4\quad 安全裕度分位点对问题二全年费用的影响（单位：万元）}{cccc}
{分位点 $q$ & 计划购电费 & 紧急购电费 & 总费用}
{不加裕度 & 1111.3 & 1284.7 & 2396.0 \\
 0.50 & 1382.0 & 446.0 & 1828.0 \\
 0.70 & 1504.9 & 211.4 & 1716.3 \\
 \textbf{0.80} & \textbf{1570.3} & \textbf{132.4} & \textbf{1702.6} \\
 0.90 & 1657.0 & 68.8 & 1725.8}

\begin{figure}[H]
  \centering
  \includegraphics[width=0.72\textwidth]{margin_tradeoff.pdf}
  \caption{安全裕度分位点与全年费用的关系，最低点在 $q=0.80$}
  \label{fig:margin}
\end{figure}

图~\ref{fig:margin} 显示的 U 形曲线正是报童模型的典型形态：分位点越低，紧急购电费越高；分位点越高，计划购电费越高；两者之和在 $q=0.80$ 处最小。

\subsection{全年结果}

2025 年 2 月 1 日至 12 月 31 日共 334 天，结果汇总于表 5。

\threelinetable{表 5\quad 问题二全年结果汇总}{lr}
{指标 & 数值}
{实际负载合计 / kWh & 37008079.49 \\
 计划购电量 / kWh & 24368275.98 \\
 计划购电费 / 元 & 15702558.44 \\
 紧急购电量 / kWh & 343800.18（占负载 $0.929\%$） \\
 紧急购电费 / 元 & 1323841.81 \\
 \textbf{总购电费用 / 元} & \textbf{17026400.25}}

全年 334 天中有 248 天出现过紧急购电，但单次规模有限：逐日紧急购电量的中位数为 $561.19\,\mathrm{kWh}$，最大的一天为 $14118.4\,\mathrm{kWh}$，说明安全裕度已经把大部分预测误差消化在计划内。

\threelinetable{表 6\quad 问题二指定日期的计划购电与紧急购电}{lrrr}
{日期 & 计划购电量 / kWh & 紧急购电量 / kWh & 费用 / 元}
{2025-03-20 & 65644.45 & 1074.66 & 45119.57 \\
 2025-06-21 & 79792.46 & 0.00 & 51753.45 \\
 2025-09-23 & 66247.34 & 1593.46 & 50655.49 \\
 2025-12-21 & 95755.68 & 649.18 & 63386.35}

\begin{figure}[H]
  \centering
  \includegraphics[width=0.95\textwidth]{display_days.pdf}
  \caption{四个指定日期的购电量与储电量，以及问题三调整后的对比}
  \label{fig:days}
\end{figure}

2025 年 6 月 21 日是夏至，光伏出力最强，虽然计划购电量不是全年最大，但当天紧急购电为 0；12 月 21 日负载最重、光伏最弱，计划购电量达 $95755.68\,\mathrm{kWh}$，为四个展示日期中最大，费用也最高。四个日期的紧急购电都集中在夜间与清晨的负载爬升段（如 3 月 20 日的 0:20--0:40、12 月 21 日的 3:30--3:50），因为该时段既没有光伏支撑、又是预测误差最容易被放大的时段。
